% Shared style for the L11a profit diagrams, \input by each Fig-L11a-*.tex.
%
% Drawn on October 23, 2026 in the style of the CHEME 5800 L9c convexity and
% gradient-descent schematics: arrowed axes, a thick ink curve (the profit), a
% blue companion curve (the payoff), open circle markers, and dashed drops to
% colored axis labels. Each figure has two panels; each panel is a schematic,
% so the axes carry no numbers.
%
% Panel coordinates: x is the share price at expiration S(T), with the y-axis
% at S(T) = 0. Inside \begin{profitdata} ... \end{profitdata} the y coordinate
% is the profit (USD/share): the scope shifts the zero-profit line up to
% \yzero and scales profits by \yscale, so every path is written in the
% structure's own units. Nodes are not scaled.
%
% Math uses the standard TeX fonts (no unicode-math), as in the L9c figures,
% so S, K, and the double dagger match the notebook's math.
%
% The notebook SVG carries a dark palette added by theme_svg.py; keep new
% colors in that script's mapping when changing the palette here.
\usepackage[no-math]{fontspec}
\IfFontExistsTF{Arial}{\setsansfont{Arial}}{\setsansfont{TeX Gyre Heros}}
\renewcommand{\familydefault}{\sfdefault}
\usepackage{amsmath,amssymb}
\usetikzlibrary{arrows.meta,calc}

% Same colors as the L9c figures.
\definecolor{vnink}{HTML}{202124}
\definecolor{gdblue}{HTML}{0091FF}
\definecolor{gdred}{HTML}{E02020}
\definecolor{gdgray}{HTML}{6D7278}
% The L9c simulated-annealing orange, used only for profits before expiration.
\definecolor{saorange}{HTML}{F08C00}

% Panel geometry (x units): the axes run to \xaxis and \yaxis, the zero-profit
% line sits at \yzero, and a profit of 1 USD/share is \yscale units tall.
\def\xaxis{9.6}
\def\yaxis{6.3}
\def\yzero{3.9}
\def\yscale{0.8}
\def\panelgap{10.6}

\tikzset{
  every picture/.append style={x=0.8cm, y=0.8cm, font=\sffamily\small, text=vnink},
  axis/.style={draw=vnink, line width=0.8pt, -{Latex[length=2.6mm, width=1.8mm]}},
  zero/.style={draw=gdgray, line width=0.5pt},
  profit/.style={draw=vnink, line width=1.6pt, line join=round},
  payoff/.style={draw=gdblue, line width=1.2pt, line join=round},
  more/.style={-{Latex[length=2.4mm, width=1.7mm]}},
  drop/.style={draw=gdgray, line width=0.8pt, dash pattern=on 3pt off 2pt},
  gap/.style={draw=vnink, line width=0.8pt,
    {Latex[length=1.8mm, width=1.3mm]}-{Latex[length=1.8mm, width=1.3mm]}},
  pt/.style={circle, fill=white, line width=1.0pt, inner sep=0pt, minimum size=2.4mm},
  note/.style={text=gdgray, font=\sffamily\footnotesize, align=center},
}

% Axes, zero-profit line, the bold run-in title (#1), and the axis labels
% (#2 horizontal, #3 vertical).
\newcommand{\panelaxes}[3]{%
  \draw[zero] (0, \yzero) -- (\xaxis - 0.2, \yzero);
  \draw[axis] (-0.15, 0) -- (\xaxis, 0);
  \draw[axis] (0, -0.15) -- (0, \yaxis);
  \node[anchor=east, text=gdgray] at (-0.05, \yzero) {$0$};
  \node[anchor=north] at (\xaxis/2, -1.05) {#2};
  \node[rotate=90, anchor=south] at (-0.55, \yaxis/2) {#3};
  \node[anchor=south] at (\xaxis/2, \yaxis + 0.15) {#1};
}
% The expiration diagrams' frame.
\newcommand{\panelframe}[1]{%
  \panelaxes{#1}{Share price at expiration $S(T)$}{Payoff and profit (USD/share)}}

% Scope in which (S, profit) are the coordinates.
\newenvironment{profitdata}{\begin{scope}[shift={(0, \yzero)}, yscale=\yscale]}{\end{scope}}

% A strike: dashed drop from the profit kink (#1, #2) to the S(T) axis, with the
% strike symbol #3 and, below it, the leg written at that strike, #4.
\newcommand{\strike}[4]{%
  \draw[drop] (#1, #2) -- (#1, -\yzero/\yscale);
  \node[anchor=north, align=center, inner sep=1.5pt, text=gdgray]
    at (#1, -\yzero/\yscale - 0.08/\yscale)
    {#3\\[-1.5pt]{\footnotesize\color{vnink}#4}};}

% The share price when the trade opens, S(0): a short tick on the S(T) axis at
% #1, labeled just above the axis so it clears the strike labels below it. Added
% October 27, 2026, so students can see where the strikes sit relative to today's
% price (the time-decay figure marks S(0) the same way).
\newcommand{\entryprice}[1]{%
  \draw[draw=vnink, line width=0.8pt] (#1, -\yzero/\yscale) -- (#1, -\yzero/\yscale + 0.2/\yscale);
  \node[anchor=south, inner sep=1.5pt, text=gdgray] at (#1, -\yzero/\yscale + 0.2/\yscale) {$S(0)$};}

% A breakeven: open red circle on the zero line at #1, label #2 placed #3 of
% the circle (above left or above right), clear of the circle's outline.
\newcommand{\breakeven}[3]{%
  \node[pt, draw=gdred,
    label={[text=gdred, inner sep=0.5pt, label distance=2pt]#3:#2}] at (#1, 0) {};}

% Maximum profit: arrow at S = #1 from the payoff (zero) up to the profit #2.
% For a credit structure the gap is the credit, grown to expiration.
\newcommand{\maxprofit}[2]{%
  \draw[gap] (#1, 0.06) -- (#1, #2 - 0.06);
  \node[anchor=south, inner sep=4pt] at (#1, #2) {max profit $=$ credit};}

% Maximum loss: arrow at S = #1 from zero down to the profit #2 (negative),
% label anchored #3 (west puts it to the right of the arrow).
\newcommand{\maxloss}[3]{%
  \draw[gap] (#1, -0.06) -- (#1, #2 + 0.06);
  \node[anchor=#3, inner sep=4pt] at (#1, #2/2) {max loss};}

% The structure's outlook (#1), as in the lecture's composite-contract table, in
% the panel's upper-right corner, the one corner no structure draws into. Added
% October 28, 2026, at the instructor's request.
\newcommand{\outlook}[1]{%
  \node[anchor=north east, inner sep=2pt, font=\sffamily\footnotesize]
    at (\xaxis, \yaxis) {{\color{gdgray}Outlook:} #1};}

% A loss that keeps growing as S(T) rises: note with its lower-left corner at
% (#1, #2), placed above the descending profit line near its arrowhead.
\newcommand{\unlimited}[2]{%
  \node[note, anchor=south west, inner sep=0pt] at (#1, #2) {unlimited\\loss};}
